\section{کاربردها و حدود انطباق مدل}

مدل این گزارش یک مدل انتزاعی با سه فرض محدودکننده است: ظرفیت هر سطل یک‌بعدی و یکسان است، روابط چسبندگی قیود سخت و از پیش معلوم‌اند، و تابع هدف فقط تعداد سطل‌ها را کمینه می‌کند. ازاین‌رو، موارد زیر زمینه‌های قابل دفاع برای استفاده از این انتزاع‌اند، نه شواهدی بر انطباق دقیق همه جزئیات عملی با مدل.

\subsection{برش یک‌بعدی مواد (\lr{Cutting Stock})}\label{subsec:lr_cutting_stock}
مسئله برش یک‌بعدی در صنایعی مانند فولاد، کاغذ، شیشه، مبلمان و نساجی رخ می‌دهد و از نظر تخصیص طول‌های موردنیاز به مواد اولیه با ظرفیت محدود، با بسته‌بندی یک‌بعدی ارتباط ساختاری دارد. تنیر و همکاران گونه‌ای صنعتی را بررسی می‌کنند که در آن اقلام می‌توانند تقسیم و سپس با جوش‌کاری دوباره متصل شوند و هدف، هم‌زمان شامل ضایعات و تعداد جوش‌ها است \cite{D1_Tanir2016}. این گونه با مسئله حاضر یکسان نیست: تقسیم‌پذیری و تابع هدف چندمعیاره در مدل ما وجود ندارد و آن مقاله نیز قید \lr{must-link} این گزارش را معرفی نمی‌کند. تنها اگر در یک کاربرد برش، چند قطعه الزاماً از یک شاخه ماده اولیه بریده شوند و سایر جزئیات نادیده گرفته شوند، می‌توان آن گروه را به‌صورت یک مؤلفه چسبندگی مدل کرد.

\subsection{زمان‌بندی کانتینرها و گردش‌کارها (\lr{Cloud Orchestration})}\label{subsec:lr_cloud_orchestration}
در سامانه‌های کانتینری، نگاشت وظایف به منابع پردازشی و رعایت ترتیب اجرای گردش‌کار از مسائل اصلی زمان‌بندی است. چارچوب \lr{KubeAdaptor} برای اجرای گردش‌کار روی کوبرنتیز، سازگاری ترتیب زمان‌بندی، مصرف منابع و اشتراک داده از طریق حجم‌های پویا را بررسی می‌کند \cite{C1_Shan2022}. مرور زمان‌بندی کربن‌آگاه نیز نشان می‌دهد که زمان‌بندی کانتینرها می‌تواند اهدافی مانند بهره‌وری منابع، زمان پاسخ، انرژی و انتشار کربن داشته باشد \cite{C4_Yang2025}.

هیچ‌یک از این دو منبع، مسئله خود را دقیقاً به‌صورت \lr{BPP} یک‌بعدی با قیود هم‌مکانی سخت فرمول‌بندی نمی‌کند. نگاشت به مدل حاضر فقط تحت یک ساده‌سازی مشخص معتبر است: هر وظیفه یک تقاضای اسکالر دارد، نودها ظرفیت یکسان دارند، هر گروه تعیین‌شده باید روی یک نود واحد اجرا شود، و هدف کمینه‌کردن تعداد نودهای فعال است. در سامانه واقعی، ظرفیت‌های چندبعدی پردازنده و حافظه، ناهمگنی نودها، زمان، تقدم وظایف، دسترس‌پذیری و اهداف انرژی می‌توانند این هم‌ارزی را از بین ببرند. بنابراین این حوزه یک انگیزه مدل‌سازی مشروط است، نه شاهدی برای ادعای تطابق دقیق یا کاهش تضمین‌شده تأخیر و کربن به‌وسیله مدل حاضر.
